\chapter{On accuracy and stability}

So far we have discussed various practical aspects of the forward and backward 
Euler methods and applied them to solve systems of ODEs. However, we 
have noticed various pecular things about these algorithms. 
For example, in \Cref{fig:test_ode_forward_backward_euler_comparison} we 
saw that the numerical solutions deviated from the true solution of 
the test ODE. In fact, in some regimes the forward Euler solution 
became unstable while the backward method remained stable. 
We will now dig deeper into these ideas starting with 
the concept of numerical accuracy. 

\section{Numerical accuracy}

When working with ODEs we typically consider two types of errors: 
a global error and a local truncation error.
Let us consider scalar ODEs. Then the global error is 
simply the difference between the true solution and 
the numerical solution at a given time step, i.e., 
\begin{equation}\label{def:global_error}
    e_n = u(t_n) - u_n, \qquad n = 1, 2, \dots,
\end{equation}
Naturally you would consider $|e_n|$ as a measure of the error at each time step
since the sign of the error is often less important than its magnitude.
The truncation error, on the other hand, measures the error 
introduced by the numerical scheme in a single step. For example, 
let us consider the forward Euler method. Then the local truncation 
error at step $n$ takes the form 
\begin{equation}\label{def:local_truncation_error}
    \tau_n = \frac{u(t_{n}) - u(t_{n-1})}{\Delta t} - f(u(t_{n-1})) 
    =: \mcl{N}_{\Delta t}[u(t_n)], 
    \qquad n =  1, 2, \dots,
\end{equation}
that is, the residual of the discretized ODE  after 
plugging in the true solution $u(t_n)$. 


\subsection{Measuring accuracy and its order}

\subsection{Modifying forward Euler to obtain second order accuracy}

\section{Numerical stability}

\subsection{Definition of stability}



\bibliographystyle{plainnat}
\bibliography{references}